% Section 4: what the certified head buys. On the engine's states: setup and the
% fixed-configuration rule, recomputation, fallback and agreement, the engine check, head-only
% runtime. Served: the stock baseline (arms, sessions, c, x, y, the envelope, Figure
% fig:frontier, the exactness class, GSM8K), then the certified head. Engine changes outside the
% head and Table tab:levers: Section 6 (levers.tex).
% Claims and status: pins from SETUP.md and the evidence READMEs; capture, alignment, stock
% divergences, held-out replay, kernel replay, engine check and head-only microbenchmark from
% committed evidence; bench's confirmation frontier, equality and quality checks; three
% sessions of the certified head served (cert-served); the BF16 paths (bf16-paths).
% Figure fig:rho (Section 5's) is placed here so that it floats near Section 5's start.
\section{What the Certified Head Buys}\label{sec:measure}
We measure the certified head on states recorded from the engine, then serve it. Its numbers come from six sources, each named with its decision rule where it is used (Table~\ref{tab:sources} in Appendix~\ref{app:register}): the \emph{held-out replay} in FP64, the \emph{kernel replay} at the captured batch shapes, the \emph{engine check} inside SGLang, a head-only \emph{microbenchmark}, the \emph{speculative captures} and the served sweeps.

\begin{figure*}[t]
\centering
% Figure fig:candidates. (a) reads evidence/head_geometry/selfevidence_plain4b_ccdf.csv.
% (b): the curves are derived, 1-(1-p)^B with the per-row interval-form shares of the
% held-out rows in evidence/head_geometry/rstock_plain4b.json (rules/gamma_1.19e-4/bucket/
% positions_needing_stock_kernel = 0.0034970857618651123 and rules/tensor_core_model/
% bucket/positions_needing_stock_kernel = 0.013988343047460449; these exist only in JSON,
% so they are copied here at full precision, and the legend's p and the shares at B = 128
% are computed from them). The markers are measured: the GPU kernel's share of batches of M
% consecutive rows of the first 60,000 capture rows that need a fallback
% (evidence/certified_head/head_path_time.csv, batch_fallback_rate and
% hopper_batch_fallback_rate).
% (c) reads evidence/certified_head/head_path_time.csv: the W8A16 path's time per call over
% the stock head and argmax (stock_us) on a decided batch (w8a16_us), and the expected time
% with the whole-batch fallback under each model (w8a16_expected_us,
% hopper_w8a16_expected_us) and with the column fallback under the conservative model
% (columns_expected_us). A microbenchmark of the head alone, probes on.
\pgfmathsetmacro{\phop}{0.0034970857618651123}
\pgfmathsetmacro{\pcons}{0.013988343047460449}
\pgfmathsetmacro{\fbhop}{100*(1-(1-\phop)^128)}
\pgfmathsetmacro{\fbcons}{100*(1-(1-\pcons)^128)}
\begin{tikzpicture}
\begin{groupplot}[group style={group size=3 by 1,horizontal sep=1.25cm},
  width=5.9cm,height=3.55cm,grid=major]
\nextgroupplot[xmode=log,ymode=log,xmin=1,xmax=248320,ymin=1e-4,ymax=1.5,
  xlabel={Candidate rows $k$},ylabel={Positions needing $\ge k$ rows},
  title={\paneltitle{a}{Candidates (measured)}},unbounded coords=discard,
  xtick={1,10,100,1000,10000,100000,248320},
  xticklabels={1,10,$10^2$,$10^3$,$10^4$,$10^5$,$V$},
  ytick={1,0.1,0.01,0.001,0.0001},yticklabels={100\%,10\%,1\%,0.1\%,0.01\%},
  minor tick num=0,xminorticks=false,yminorticks=false,
  legend style={at={(0.97,0.03)},anchor=south east},
  legend image post style={xscale=0.5}]
\addplot[series int8 row] table[x=k,y=plain_int8_row_greedy,col sep=comma]{../evidence/head_geometry/selfevidence_plain4b_ccdf.csv};\addlegendentry{int8 per row}
\addplot[series int8 g128] table[x=k,y=plain_int8_g128_greedy,col sep=comma]{../evidence/head_geometry/selfevidence_plain4b_ccdf.csv};\addlegendentry{int8 128 cols}
\addplot[vpblue,semithick,only marks,mark=o,mark size=1.4pt] table[x=k,y=plain_int8_g128_gumbel_t1.0,col sep=comma]{../evidence/head_geometry/selfevidence_plain4b_ccdf.csv};\addlegendentry{same, $T{=}1$}
\addplot[series fp8] table[x=k,y=plain_fp8_row_greedy,col sep=comma]{../evidence/head_geometry/selfevidence_plain4b_ccdf.csv};\addlegendentry{FP8 per row}
\addplot[series int4] table[x=k,y=plain_int4_g128_greedy,col sep=comma]{../evidence/head_geometry/selfevidence_plain4b_ccdf.csv};\addlegendentry{int4 128 cols}
\nextgroupplot[xmode=log,log basis x=2,xmin=1,xmax=256,ymin=0,ymax=1,
  xtick={1,4,16,64,256},xticklabels={1,4,16,64,256},
  xlabel={Rows in the batch $B$},ylabel={Batches with a fallback},
  ytick={0,0.2,0.4,0.6,0.8,1},yticklabels={0\%,20\%,40\%,60\%,80\%,100\%},
  title={\paneltitle{b}{Whole-batch fallback}},
  legend style={at={(0.03,0.97)},anchor=north west}]
\addplot[series conservative,domain=1:256,samples=120]{1-(1-\pcons)^x};\addlegendentry{cons., derived}
\addplot[series hopper,domain=1:256,samples=120]{1-(1-\phop)^x};\addlegendentry{Hopper, derived}
\addplot[only marks,mark=square*,mark size=1.5pt,vpgreen] table[x=M,y=batch_fallback_rate,col sep=comma]{../evidence/certified_head/head_path_time.csv};\addlegendentry{cons., kernel}
\addplot[only marks,mark=*,mark size=1.6pt,ink] table[x=M,y=hopper_batch_fallback_rate,col sep=comma]{../evidence/certified_head/head_path_time.csv};\addlegendentry{Hopper, kernel}
\draw[muted,densely dotted] (axis cs:128,0) -- (axis cs:128,1);
\node[annot,anchor=south east] at (axis cs:124,{\fbcons/100}) {\vppct{0}{\fbcons}};
\node[annot,anchor=south east] at (axis cs:122,{\fbhop/100+0.01}) {\vppct{0}{\fbhop}};
\nextgroupplot[xmode=log,log basis x=2,xmin=1,xmax=256,ymin=0.5,ymax=2.5,
  xtick={1,4,16,64,256},xticklabels={1,4,16,64,256},
  ytick={0.5,1,1.5,2,2.5},yticklabels={0.5,1,1.5,2,2.5},
  xlabel={Rows in the batch $M$},ylabel={Time over the stock head},
  title={\paneltitle{c}{Head-path time (measured)}},
  legend style={at={(0.03,0.97)},anchor=north west},
  every axis plot/.append style={mark size=1.4pt}]
\draw[muted,densely dotted] (axis cs:1,1) -- (axis cs:256,1);
\addplot[series int8 row,mark=*] table[x=M,y expr={\thisrow{w8a16_us}/\thisrow{stock_us}},col sep=comma]{../evidence/certified_head/head_path_time.csv};\addlegendentry{decided batch}
\addplot[series conservative,mark=square*,mark options={solid}] table[x=M,y expr={\thisrow{w8a16_expected_us}/\thisrow{stock_us}},col sep=comma]{../evidence/certified_head/head_path_time.csv};\addlegendentry{whole batch, cons.}
\addplot[series hopper,mark=triangle*] table[x=M,y expr={\thisrow{hopper_w8a16_expected_us}/\thisrow{stock_us}},col sep=comma]{../evidence/certified_head/head_path_time.csv};\addlegendentry{whole batch, Hopper}
\addplot[vpgreen,thick,dashdotted,mark=diamond*,mark options={solid}] table[x=M,y expr={\thisrow{columns_expected_us}/\thisrow{stock_us}},col sep=comma]{../evidence/certified_head/head_path_time.csv};\addlegendentry{columns, cons.}
\end{groupplot}
\end{tikzpicture}
%
\caption{Int8 leaves few candidates, but whole-batch fallback grows with the batch; timed alone, with the per-position fallback under the \modelcons{}, the head is expected to take about 0.7 of the stock head's time up to 16 positions, and every mode is slower from 128. \textbf{(a)}~Held-out replay, \Rreal{}: four copies, greedy, and the per-128-weights int8 copy in a seeded race at $T=1$ \ev{ccdf}. \textbf{(b)}~Lines: independent fallback at Table~\ref{tab:fallback}'s held-out shares \ev{rstock}; markers: kernel replay \ev{kernel-path}. \textbf{(c)}~Over the stock head: measured with every position certified (medians; p10--p90 within the markers below 256), and expected with each fallback, from measured times and (b)'s conservative rates \ev{kernel-path}. Shaded: consecutive replay positions, not engine batches.}
\label{fig:candidates}\label{fig:geometry-candidates}
\end{figure*}

\subsection{Setup}\label{sec:setup}\label{sec:transport-capture}
\paragraph{System} One NVIDIA GH200 with SGLang~\cite{sglangpaper} at commit \sourcepin{}, PyTorch 2.13 for CUDA 13.0 and, unless a configuration names another, the FlashInfer attention backend (Appendix~\ref{app:pins} lists every version). The target is Qwen3.5-4B at revision \modelpin~\cite{qwencard}, whose decoder layers are, three in four, Gated DeltaNet (GDN) layers: linear attention with a recurrent state per request~\cite{gdn}. The drafters are MTP-4B, the model's native MTP layer, run with top-$k$ 1 and four draft steps in the captures (three in Figure~\ref{fig:breakdown} and served), and DFlash-4B, the public block drafter for this target~\cite{dflashpaper,dflash4bcard}, which drafts blocks of 16 tokens (8 in one tuned arm) and projects them through the target's head.

\paragraph{Capture} A capture-only engine patch records each position's head input as BF16 bit patterns, with the token the engine chose; under speculation it also records the drafter's head inputs and pairs each drafted token with the target head input that verifies it. The capture servers ran without CUDA graphs or the overlap scheduler, with at most 16 running requests in plain decoding (Appendix~\ref{app:transport-capture}). The 320 public \emph{capture prompts} were decoded greedily for 384 tokens and split by prompt into an analysis half, on which the data-dependent choices were made, and a held-out half. The deployed copy and its bound depend on the weights alone, so the kernel replay draws positions from both halves. In plain decoding the FP64 argmax of the captured head input equals the engine's token at 99.51\% of positions, and every disagreement lies within one BF16 spacing of the FP64 winner \ev{alignment}.
\paragraph{A fixed configuration}\label{sec:same-shape} Bit-equality is a useful test only within one configuration, because the stock engine's own output varies across configurations. On a server that runs at most 16 requests at once and queues the rest, plain greedy decoding differs within 256 tokens on 52\% of 320 separate prompts between a client that keeps 1 request in flight and one that keeps 32 (Appendix~\ref{app:divergence-rates}) \ev{state-noise}. A bit-exact trace places the first differences upstream of the head, in kernels whose results depend on the batch or the memory pools, never in the head GEMM or the tie rule \ev{state-mech}. Every comparison therefore fixes the serving configuration and compares speculative runs with stock speculative verification.

\subsection{How Many Entries Are Recomputed}\label{sec:geometry-selfevidence}\label{sec:candidates}
\paragraph{Kernel replay, \Rstock{}} On the first 60,000 captured positions, the certified head recomputes about 1.8 entries per decision on average and one at the median, depending on the stock-kernel model through the margin $\delta_{\max}$ (Table~\ref{tab:fallback} in Appendix~\ref{app:candidates}) \ev{kernel-replay}.

\paragraph{Held-out replay, \Rreal{}} This replay compares copies of the head on 6,005 held-out plain-decoding positions, drawn as whole decode steps, with the simpler \Rreal{} decision, whose screen needs no margin. No error bound was violated, and the exact winner was always a candidate \ev{selfevidence,selfevidence-csv}. Int8 with one scale per entry needs 1.55 candidates on average; FP8 needs about three times as many, and int4 leaves tens of thousands (Figure~\ref{fig:candidates}a; Table~\ref{tab:candidates}). Int8 is enough because the margins are wide: its intervals are about two logits wide, while the median gap between the two largest logits is 7.3 logits \ev{stats}, so a tighter bound would help little (Appendix~\ref{app:candidates}).

\subsection{Fallback and Agreement}\label{sec:rstock-share}
\paragraph{Kernel replay, \Rstock{}, interval condition} The certified head leaves 0.28\% of positions to the stock head under the \modelhopper{} and 1.46\% under the \modelcons, and with that fallback every token equals the engine's \ev{kernel-replay}. The held-out replay gives nearly the same shares, 0.35\% and 1.40\%; under the gap condition, which needs a margin of more than one BF16 spacing, its share never falls below about 2\% (Table~\ref{tab:fallback}). This agreement is consistent with each model, not a proof of any: even an IEEE FP32 summation tree, with a $\gamma_{\rm stock}$ about 360 times smaller than the conservative one, certifies every held-out position and gets every token right \ev{rstock}.

Fallback grows with the batch, because one undecided position sends the whole batch to the stock head. If positions fell back independently with probability $p$, a call of $M$ positions would contain an undecided one with probability $1-(1-p)^M$, which matches the measured share of decode steps with a fallback (Table~\ref{tab:fallback}). On calls of 128 consecutive positions the kernel replay measures 27\% and 77\% (Hopper, conservative), where the held-out rates predict 36\% and 84\% (Figure~\ref{fig:candidates}b; Appendix~\ref{app:candidates}) \ev{kernel-path}. At serving batch sizes the certified head therefore depends on the per-position fallback.

\paragraph{Engine check, \Rstock{}, interval condition}\label{sec:engine} The engine check ran the kernel inside SGLang in \emph{check mode}, which also runs the stock head and counts the positions whose token differs, untimed and under the \modelcons, on plain decoding, greedy and fixed-noise verification and the drafters' projections. In concurrent runs (up to 16 requests for plain decoding, 8 for speculation) none of 196,734 certified positions differed from the stock head at the same batch shape \ev{kernel-engine}. Runs with one request at a time gave identical certified and stock outputs on all 64 prompts for plain decoding and MTP-4B; DFlash-4B's were compared only when served (Section~\ref{sec:served}). Every path used the whole-batch fallback except one plain-decoding run, which used the per-position one. Plain decoding fell back at 1.38\% of positions (kernel replay: 1.46\%) and MTP-4B verification at 1.74\%. DFlash-4B verification sends up to 128 positions per call and left only 3.99\% of them open, yet ran the whole-batch fallback in 90\% of its verify head calls.

\subsection{Head-Only Runtime}\label{sec:kernel-runtime}
A head call saves time only if
\begin{equation}
T_{\rm pass}+T_{\rm select}+T_{\rm rescore}+P_{\rm fb}(M)\,T_{\rm fallback}(M)<T_{\rm head}(M),\label{eq:headgate}
\end{equation}
where the left side adds the int8 pass, the selection of candidates, their recomputation and the fallback's time weighted by the share $P_{\rm fb}(M)$ of calls that fall back, and $T_{\rm head}$ is the stock head (GEMM, FP32 copy and argmax).

In the microbenchmark (the head alone at $M=1$ to 256, CUDA graphs, warm L2 cache; medians of 30 graph replays) a fully certified batch at $M=1$ takes 239.6~\textmu s, against 370.5~\textmu s for the stock head. Adding each fallback's measured cost at its measured rate gives an expected time of about 0.7 of the stock head's up to 16 positions, with the per-position fallback under the \modelcons{} (Figure~\ref{fig:candidates}c) \ev{kernel-rowwise,kernel-path}. From 128 positions every mode is slower, because the int8 pass alone is slower than the stock head there; with the whole-batch fallback, 77\% of calls also pay for the stock head. The pass is not bound by memory bandwidth: at one position its directed-rounding epilogue keeps SM throughput at 79\% of peak while it draws 62--63\% of DRAM bandwidth, so its time grows with $M$ faster than the stock GEMM's \ev{kernel-micro}. At one position the expected time under the \modelcons{} is 0.66--0.67 of the stock head's with either fallback \ev{kernel-path}, so with the head at 10.3\% of a plain decode step at batch 1, the step would shorten by at most 3.4--3.5\% (derived; served, 2.5\%, Section~\ref{sec:served}).

\subsection{Served Decoding}\label{sec:served}
\paragraph{The stock baseline} Each serving configuration, or \emph{arm}, runs with CUDA graphs and the overlap scheduler (checked at launch) and without the prefix cache, which the workload does not reuse. A \emph{session} launches each compared arm afresh, once, at full capacity, and sweeps client concurrency $c$, the requests a closed-loop client keeps in flight, on 1,152 held-out public prompts (chat, code, mathematics) with 512-token outputs \ev{workload}. Figure~\ref{fig:frontier} plots the per-request output rate $x$, including the time to first token, against throughput $y$ (Appendix~\ref{app:serving}); the \emph{envelope} is the best arm at each concurrency. Over three sessions the best arm of each drafter serves more than plain decoding up to $c=32$ and, with SGLang's default admission, less from 48 \ev{bench-frontier}. At $c=1$, DFlash-4B with 16-token blocks gives 3.06 times plain decoding's throughput, both with Triton attention. Every arm is stock or \emph{exact up to rounding}, a class fixed after the first comparisons had been seen: each first divergence from a matched stock reference is an exact tie, one BF16 step or a near tie within two, and no output differs in length \ev{bench-equality}. On GSM8K the four arms tested (two MTP, 8-token DFlash, and Section~\ref{sec:levers}'s replayed plain decoding) are within plain decoding's noise, which excludes losses above about 3.3 points (derived), not smaller ones \ev{bench-quality}.

\begin{figure}[t]
\centering
% Figure fig:frontier (Section 4.5): the confirmed serving frontier. Its width is the line width up
% to 10.5 cm and its height 0.48 of that. Reads
% evidence/bench/confirm/frontier.csv (per arm and client concurrency: means over sessions
% of x = output tokens per second per request, end to end, and y = output tokens per second
% on the GPU; n = valid sessions) and evidence/bench/confirm/envelope-all.dat (the Pareto
% envelope over the points with at least three valid sessions, grey band). Points with fewer
% than three sessions are hollow; they do not rank on the frontier (evidence/bench/README.md).
% Styles follow style.tex, one colour per arm: plain decoding (arm plain) ink, solid, circles;
% DFlash with 8-token blocks (arm dflash8) violet, solid, diamonds; DFlash with 16-token blocks
% (arm dflash16) reddish purple, dashed, diamonds; MTP (arm mtp) green, dash-dot, triangles,
% filled where three or more sessions are valid and hollow (mark hollow) where only one is, which
% the caption says. Each curve runs from c = 1 (lower right) to c = 128 (upper left); the plain
% curve's ends are labelled, clear of its line (c = 128 above and left of its point, clear of MTP's c = 128
% triangle below it and of the top frame). Error bars are one standard deviation over sessions (x_e2e_std, y_std); every
% one is below 2% of its mean and hidden by its marker, which the caption says. No arm uses the
% certified head. The figure has one panel, so it has no title inside it; the caption names it
% and its status.
% Direct labels, no legend, each placed where its curve is alone (positions checked against
% frontier.csv): plain decoding left of its c = 1-2 run (nothing lies left of x = 228 below
% y = 3,600), opposite the c = 1 label; MTP, two lines, right of its c = 1 point (461, 456), under the DFlash c = 1
% points; 8-token DFlash above and right of its c = 16 point (370, 5,256), where the envelope is
% the upper boundary of every curve; 16-token DFlash under its c = 48 point (130, 5,326), the
% lowest curve left of x = 184, clear of the plain curve at x >= 228.
% Keep only points with at least #1 sessions, or with fewer than #1.
\pgfplotsset{
  only if n at least/.style={x filter/.append code={%
    \ifnum\thisrow{n}<#1\relax\def\pgfmathresult{inf}\fi}},
  only if n below/.style={x filter/.append code={%
    \ifnum\thisrow{n}<#1\relax\else\def\pgfmathresult{inf}\fi}}}
\pgfmathsetmacro{\frontierw}{min(\linewidth,10.5cm)}
\pgfmathsetmacro{\frontierh}{0.48*\frontierw}
\begin{tikzpicture}
\begin{axis}[xmode=log,ymode=log,width=\frontierw pt,height=\frontierh pt,
  xmin=65,xmax=1150,ymin=180,ymax=25000,grid=major,
  xlabel={Per-request rate $x$ (output tokens/s)},
  ylabel={Throughput $y$ (output tokens/s)},
  xtick={100,200,500,1000},xticklabels={100,200,500,1000},
  ytick={300,1000,3000,10000},yticklabels={300,1k,3k,10k},
  % one standard deviation over sessions; keep last in a plot's options (it changes key path)
  sd bars/.style={error bars/.cd,x dir=both,x explicit,y dir=both,y explicit,error bar style={thin}}]
\addplot[line width=4pt,ccontext!25,forget plot,no markers]
  table[x=x,y=y]{../evidence/bench/confirm/envelope-all.dat};
\addplot[arm plain,discard if not={label}{plain-tuned},sd bars]
  table[x=x_e2e_mean,y=y_mean,x error=x_e2e_std,y error=y_std,col sep=comma]{../evidence/bench/confirm/frontier.csv};
\addplot[arm dflash8,discard if not={label}{dflash-tuned},sd bars]
  table[x=x_e2e_mean,y=y_mean,x error=x_e2e_std,y error=y_std,col sep=comma]{../evidence/bench/confirm/frontier.csv};
\addplot[arm dflash16,discard if not={label}{dflash-tuned-b16},sd bars]
  table[x=x_e2e_mean,y=y_mean,x error=x_e2e_std,y error=y_std,col sep=comma]{../evidence/bench/confirm/frontier.csv};
\addplot[arm mtp,discard if not={label}{mtp-tuned},sd bars]
  table[x=x_e2e_mean,y=y_mean,x error=x_e2e_std,y error=y_std,col sep=comma]{../evidence/bench/confirm/frontier.csv};
% the one-session MTP points drawn again over their filled marks, hollow; the caption names them
\addplot[only marks,cmtp,mark triangle,mark hollow,forget plot,discard if not={label}{mtp-tuned},only if n below=3]
  table[x=x_e2e_mean,y=y_mean,col sep=comma]{../evidence/bench/confirm/frontier.csv};
\node[annot,anchor=west] at (axis cs:300,282) {$c{=}1$};
\node[annot,anchor=south east] at (axis cs:106,14800) {$c{=}128$};
% direct labels (positions in the comment above)
% each label in its arm's colour (style.tex), so the two DFlash labels name their lines unambiguously
\node[annot,anchor=east,text=cstock] at (axis cs:255,450) {plain decoding};
\node[annot,anchor=north west,align=left,text=cmtp] at (axis cs:478,640) {MTP, 3 steps,\\FlashInfer};
\node[annot,anchor=south west,text=cdflash8] at (axis cs:390,5700) {DFlash, 8-token blocks};
\node[annot,anchor=north,align=center,text=cdflash16] at (axis cs:135,3800) {DFlash,\\16-token blocks};
\end{axis}
\end{tikzpicture}%
%
\caption{Under SGLang's default admission, speculation leads plain decoding at low concurrency and trails it at high: the best arm of each drafter (not all drawn) serves more up to $c=32$ and less from 48. Four of nine tuned stock arms, $c=1$ to 128 (64 for 16-token DFlash); points average three or four sessions (hollow: one), with standard deviations under 2\% of each mean \ev{bench-frontier}. Grey: Pareto envelope of all nine arms' points with three or more sessions.}
\label{fig:frontier}
\end{figure}

\paragraph{The certified head, served} The four tuned arms of Table~\ref{tab:levers} were each served with and without the certified head (per-position fallback, \modelcons{}) in three sessions. The head certifies calls of at most 64 positions; a larger call runs the stock head in the same CUDA graph, the head \emph{gated off}. At $c=1$ it raises throughput by 2.5\% for plain decoding, 4.6\% for MTP and 1.2\% for 16-token DFlash, 28--71\% of the saving the head-only times predict, and keeps gaining up to $c=32$ for plain decoding and $c=8$ for MTP (Table~\ref{tab:levers}; Appendix~\ref{app:cert-served}) \ev{cert-served}. Both DFlash arms lose 1.0--3.0\% above $c=1$ where the head is active (calls of up to 28--64 positions; 54--82\% of verify calls fall back in check mode) and MTP 3.3\% at 64 (64-position draft calls, 76\% falling back), though plain decoding at $c=32$ and 64 and MTP at $c=8$--32 run such calls without a detectable loss. On the envelope it gains only at $c=1$, and loses 1.0--3.0\% at $c=2$--8 (calls of up to 30--64 positions) and 0.6\% at 128, gated off.

\begin{figure*}[t]
\centering
% Figure fig:rho. (a) is drawn to scale from medians: ||h^d|| = 50.596 and ||h^t|| = 159.2
% (evidence/head_geometry/transport_dflash4b.json, headline hdnorm2/htnorm2 medians),
% cos(h^d, h^t) = 0.408 (evidence/head_geometry/stats_dflash4b.json, pairs/cos_hd_ht/all
% q50), and the median threshold delta* = 0.00854 (evidence/precision/head_constants.json,
% tiles/contiguous/euclid_int8_row_threshold/p50). (b) reads
% evidence/head_geometry/rho_dflash4b_quantiles.csv; the MTP-4B markers are the p10,
% median and p90 of evidence/head_geometry/transport_mtp4b.json,
% headline/dnorm2_over_htnorm2 (0.7666, 0.9541, 1.1756).
\begin{tikzpicture}[x=0.0335cm,y=0.0335cm,font=\footnotesize]
\pgfmathsetmacro{\Lt}{159.2}
\pgfmathsetmacro{\Ld}{50.596}
\pgfmathsetmacro{\Cth}{0.408}
\pgfmathsetmacro{\Sth}{sqrt(1-\Cth*\Cth)}
\coordinate (O) at (0,0);
\coordinate (T) at (\Lt,0);
\coordinate (D) at ({\Ld*\Cth},{\Ld*\Sth});
\coordinate (F) at ({\Lt*\Cth*\Cth},{\Lt*\Cth*\Sth});
\path[use as bounding box] (-40,-40) rectangle (176,104);
\node[anchor=north,font=\footnotesize] at (68,104) {(a) Median DFlash-4B head inputs, to scale};
% the line through h^d and the foot of the perpendicular from h^t (best rescaling of h^d)
\draw[muted,densely dotted] (O) -- ($(O)!1.18!(F)$);
\draw[muted,densely dotted] (T) -- (F);
\node[note,anchor=west,align=left] at ($(F)+(3,4)$) {closest rescaling of $h^d$:\\$\rho\approx0.91$};
% vectors
\draw[flow,very thick,shorten >=2.2pt] (O) -- (T);
\node[below,text=ink] at (80,-2) {$h^t$, $\norm{h^t}_2\approx159$};
\draw[flow,very thick,teal] (O) -- (D);
\node[text=teal,anchor=east,align=right] at ($(D)+(-3,-10)$) {$h^d$, $\norm{h^d}_2$\\$=\sqrt{2560}$};
\draw[flow,thick,brick,dashed,shorten >=2.2pt] (D) -- (T);
\node[text=brick,anchor=south west] at (78,22) {$\Delta$: $\norm\Delta_2\approx0.92\norm{h^t}_2$};
\draw[muted] ($(O)+(13,0)$) arc[start angle=0,end angle=65.9,radius=13];
\node[note,anchor=west] at (14,5) {$\cos=0.41$};
% the ball inside which transport would beat int8: radius 0.00854*159.2 = 1.36 units, to scale
\fill[ochre] (T) circle[radius={0.00854*\Lt}];
\draw[ochre,thin] (T) circle[radius=2.6pt];
\draw[ochre,thin] (T) -- (152,-14);
\node[note,text=ochre,anchor=north,align=center] at (150,-14) {transport beats int8 only\\inside the dot in this circle:\\radius $0.0085\norm{h^t}_2$};
\end{tikzpicture}\hfill
\begin{tikzpicture}
\begin{semilogxaxis}[width=8.9cm,height=4.7cm,xmin=2e-3,xmax=1.6,ymin=0,ymax=1,
  xlabel={Drift ratio $\rho$ of a pair, or per-row threshold $\delta^\ast_i$},ylabel={Cumulative share},grid=major,
  title={(b) Measured drift against the thresholds},xminorticks=false,
  legend style={at={(0.5,-0.36)},anchor=north,legend columns=3,
    /tikz/every even column/.append style={column sep=5pt}}]
\addplot[thick,teal,dashed] table[x=threshold_int8_row,y=q,col sep=comma]{../evidence/head_geometry/rho_dflash4b_quantiles.csv};\addlegendentry{$\delta^\ast_i$, int8 per row}
\addplot[thick,teal,densely dotted] table[x=threshold_int8_g128,y=q,col sep=comma]{../evidence/head_geometry/rho_dflash4b_quantiles.csv};\addlegendentry{$\delta^\ast_i$, int8 128 cols}
\addplot[thick,ochre,dashdotted] table[x=threshold_int4_row,y=q,col sep=comma]{../evidence/head_geometry/rho_dflash4b_quantiles.csv};\addlegendentry{$\delta^\ast_i$, int4 per row}
\addplot[very thick,ink] table[x=rho_all,y=q,col sep=comma]{../evidence/head_geometry/rho_dflash4b_quantiles.csv};\addlegendentry{$\rho$, DFlash-4B pairs}
\addplot[only marks,mark=triangle*,mark size=2.6pt,brick] coordinates {(0.7666151341559915,0.1) (0.954136519993249,0.5) (1.1756316575526062,0.9)};\addlegendentry{$\rho$, MTP-4B p10, 50, 90}
\draw[{Stealth[length=1.6mm]}-{Stealth[length=1.6mm]},muted] (axis cs:0.0095,0.5) -- (axis cs:0.84,0.5);
\node[font=\scriptsize,text=ink,fill=white,inner sep=1pt] at (axis cs:0.037,0.5) {medians $\approx100\times$ apart};
\end{semilogxaxis}
\end{tikzpicture}
%
\caption{Transport fails because draft head inputs are far from the target's: the median drift is 107 (DFlash-4B) and 112 (MTP-4B) times the median threshold $\delta^\ast$ below which transport beats the per-entry int8 copy. \textbf{(a)}~Median DFlash-4B head inputs, to scale (norms over the 4,020 whole-step pairs, cosine over 20,000 random aligned pairs); transport wins only if $h^d$ lies in the dot of radius $\delta^\ast\norm{h^t}_2$ (Proposition~\ref{thm:crossover}). \textbf{(b, c)}~Drift $\rho$ over 40,000 aligned held-out pairs per drafter, against each entry's threshold $\delta^\ast_i$ from the weights; arrows join the medians \ev{head-constants,rho,transport,rho-mtp,stats-dflash}.}
\label{fig:rho}\label{fig:transport}
\end{figure*}

\paragraph{Costs and exactness} Deploying the head has further costs: the certified graphs take 3.8--25.7~GB more memory, so both DFlash comparisons pinned the attention cache at 60,000 tokens; gated off, the head still costs 0.6--1.2\% (8-token DFlash's intervals include 1); and a batch that requests log-probabilities, penalties or a grammar gets the stock verify. Its tokens equal stock's at $c=1$, in check mode at the tuned flags (0 of 592,433 certified positions differ, besides the engine check's 196,734) and, for MTP, under identical batch evolution. Closed-loop runs of plain decoding, MTP and 8-token DFlash diverge from stock at 1.03--1.21 times the rate between stock runs (only MTP's 95\% interval, 1.0002--1.22, excludes 1); 16-token DFlash's runs diverge 9 times against none (at $c=8$, gated off), all at ties or one BF16 step. The one large divergence, in a certified MTP run at $c=64$, commits the token that stock BF16 decoding at batch 1, unlike Hugging Face's, also puts first, 8.9 nats below FP32's top, at a position past the model's end of text \ev{bf16-paths}.
